\begin{result}{F.3.25}{Rational CM endomorphism orbits and genus-three normalizer recognition}{Proved; actual endomorphism and arithmetic maps}
Let $A/\Q$ be absolutely simple of dimension $g$, with
$E=\operatorname{End}^0_{\Qb}(A)$ a CM field of degree $2g$.
Then $E/\Q$ is Galois, the actual endomorphism action
$G_{\Q}\to\operatorname{Aut}_{\Q}(E)$ is surjective, and
its kernel is $G_E$. At a completely split prime its
actual action is transitive on the $2g$ CM eigenlines,
which form conjugate symplectic pairs.

For every smooth proper genus-three curve $C/\Q$ with
absolutely simple Jacobian and point $y\in C(\Q)$,
every prime $\ell$, and every finite extension $K/\Q$,
a nontrivial closed procyclic subgroup of
$\pi_1((C\setminus\{y\})_{\overline K})^{(\ell)}$
whose arithmetic normalizer has open image in $G_K$
is contained in an actual cusp inertia subgroup.
\end{result}
\begin{proof}
We first determine the exact endomorphism action from the rational
Hodge filtration, then apply the universal quotient of F.3.24.

\paragraph{The type and its right stabilizer.}
Let $L$ be the normal closure of $E$, put
$G=\operatorname{Gal}(L/\mathbf Q)$ and
$H=\operatorname{Gal}(L/E)$, and fix the inclusion
$E\hookrightarrow L\subset\overline{\mathbf Q}$.
The embeddings of $E$ are identified with the right
cosets $G/H$ by $gH\mapsto g|_E$. Let $S\subset G$
be the inverse image of the CM type occurring in
$\operatorname{Fil}^1H^1_{\mathrm{dR}}(A/\mathbf Q)
\otimes\overline{\mathbf Q}$. Thus $S$ is right
$H$-invariant and
\[
                           S\sqcup cS=G,                  \tag{CMQ.1}
\]
where $c$ is the central complex conjugation of the
Galois CM field $L$. Here this assertion about the normal
closure has an explicit field proof. The field $L$ is
the compositum of the conjugate CM fields $g(E)$.
Complex conjugation on each $g(E)$ is its intrinsic
CM involution, so every embedding carries it to the
CM involution on the image. Hence complex conjugation
commutes with $G$ on each generating subfield and
therefore on $L$. Its fixed field is totally real:
for $x$ fixed by $c$ and $g\in G$, one has
$c(gx)=g(cx)=gx$. Since $L$ contains the totally
imaginary field $E$, the involution is nontrivial.
Thus $L/L^{\langle c\rangle}$ is a totally
imaginary quadratic extension of a totally real field.

Its right stabilizer is exactly $H$. Indeed let
\[
                    K=\{k\in G:Sk=S\}.
\]
It contains $H$, and $c\notin K$ by (CMQ.1).
The fixed field $F=L^K$ is a CM subfield of $E$:
centrality of $c$ makes its restriction to $F$ a
nontrivial involution, its fixed field lies in the
totally real field $L^{\langle c\rangle}$, and no
embedding of $F$ is real because $c$ is central and
acts nontrivially. Right $K$-invariance says precisely
that the type of $E$ is the inverse image, under
$\operatorname{Hom}(E,\mathbf C)\to
 \operatorname{Hom}(F,\mathbf C)$, of a CM type of $F$.

If $K$ were strictly larger than $H$, this would give
an actual proper abelian subvariety of $A$. To see the
map explicitly, identify $H_1(A(\mathbf C),\mathbf Q)$
with the rank-one $E$-module $E$. The induced-type
condition says that its complex structure, expressed
as multiplication in $E\otimes\mathbf R$, belongs to
$F\otimes\mathbf R$. Choose an $F$-basis of $E$ and
project onto one coordinate. The resulting rational
idempotent $p:E\to E$ commutes with the complex
structure and has rational rank $[F:\mathbf Q]$.
Choose an integer $n>0$ clearing its denominators on
$H_1(A(\mathbf C),\mathbf Z)$. The map $np$ gives a
complex-linear homomorphism of the actual analytic
torus $A(\mathbf C)$. Its graph is an algebraic
subvariety of the projective product by Chow's theorem,
so it is an abelian-variety endomorphism. The base-change
map
$\operatorname{End}_{\overline{\mathbf Q}}(A)
 \to\operatorname{End}_{\mathbf C}(A_{\mathbf C})$
is an isomorphism: spread any complex endomorphism over
a finitely generated extension of $\overline{\mathbf Q}$,
which is regular, and apply the graph-descent assertion
of Milne's original Corollary~20.4(b), pp.~146--147~\cite{AuditMilne1986AV}.
This is also the rigidity argument obtained by spreading
over a connected base and using the discrete Hom functor.
The actual endomorphism $np$ therefore descends to
$\overline{\mathbf Q}$ and has image of dimension
$[F:\mathbf Q]/2$ strictly between zero and $g$.
Absolute simplicity excludes this. Hence
\[
                              K=H.                        \tag{CMQ.2}
\]

\paragraph{The semilinear de Rham equation.}
The actual action $a_\sigma$ on $E$ corresponds to
some $n_\sigma H\in N_G(H)/H$:
\[
                       a_\sigma=n_\sigma|_E.
\]
For an $E$-eigenvector $w$ of de Rham character
$\varphi$ in the rational Hodge filtration, applying
$\sigma$ coefficientwise to $e^*w=\varphi(e)w$ gives
\[
             a_\sigma(e)^*\sigma(w)
                    =\sigma(\varphi(e))\sigma(w).
\]
The new character is
$\sigma\circ\varphi\circ a_\sigma^{-1}$.
The Hodge filtration is defined over $\mathbf Q$,
so the corresponding type identity is
\[
                         \sigma S=S n_\sigma.             \tag{CMQ.3}
\]
The right action has the stated orientation because
composing $g|_E$ with $n|_E$ gives $(gn)|_E$.
This also verifies compatibility with the
homomorphism law for actual endomorphism conjugation.

If $\sigma$ fixes $L$ pointwise, the left side of
(CMQ.3) is $S$. Equations (CMQ.2)--(CMQ.3) force
$n_\sigma\in H$ and $a_\sigma=1$. Thus $a$ factors
through the finite group $G$. For this factored
homomorphism, its kernel is exactly
\[
                 H^*=\{g\in G:gS=S\},                    \tag{CMQ.4}
\]
the left stabilizer of $S$. In particular $H^*$ is
normal in $G$.

Normality turns left stabilization into right
stabilization. For $s\in S$ and $h\in H^*$, write
$sh=(shs^{-1})s$. Since $shs^{-1}\in H^*$, the
right side belongs to $S$. Hence $SH^*=S$ and
\[
                                H^*\subset H.              \tag{CMQ.5}
\]
Now the exact index comparison is
\[
 [G:H]\leq [G:H^*]=|\operatorname{im}a|
       \leq|N_G(H)/H|\leq[G:H].                            \tag{CMQ.6}
\]
Every inequality is an equality. Thus $H^*=H$,
$N_G(H)=G$, and the endomorphism action is onto
$N_G(H)/H$. The subgroup $H$ is normal, so $E/\mathbf Q$
is Galois. Since $L$ was its normal closure, $L=E$
and $H=1$. Formula (CMQ.4) therefore proves that the
kernel of the original action is exactly $G_E$.

\paragraph{The actual Tate-module permutation.}
At a completely split prime let
$V_\varphi\subset V_\ell A$ be the eigenline for
$\varphi:E\to\mathbf Q_\ell$. The arithmetic action
is $\mathbf Q_\ell$-linear, and its endomorphism
conjugation gives the exact formula
\[
 \rho_\ell(\sigma)V_\varphi
                    =V_{\varphi\circ a_\sigma^{-1}}.       \tag{CMQ.7}
\]
Surjectivity of $a$ makes the orbit all the
$2g$ embeddings. A polarization over $\mathbf Q$
has Rosati involution equal to complex conjugation
on $E$. Its alternating pairing is therefore zero
between all eigenlines except the pairs
$V_\varphi,V_{\bar\varphi}$, and is nondegenerate
on each such pair. Scaling their generators yields
an actual symplectic basis made of the orbit lines.
This supplies both the arithmetic permutation and
the pairing needed by the simultaneous normalizer
calculation.


\paragraph{The projector after the actual simultaneous constraints.}
Keep the genus-three universal submodule $U_6\subset\mathfrak h_6$
and quotient $q_6$ of F.3.24. Fix any symplectic basis of $V$.
Let $T_{\mathrm{all}}$ be the full subspace of special degree-six
derivations tangent to all six basis lines, without imposing a
character condition. Then
\[
                         q_6(U_6\cap T_{\mathrm{all}})=0.       \tag{QCM.1}
\]

Indeed, let $T_{\mathrm{diag}}$ be the full split diagonal
similitude torus for this basis, with multiplier $t$. The
subspace $T_{\mathrm{all}}$ is stable under it: for a diagonal
element $h$, conjugation sends $D(v)=[A_v,v]$ to another
bracket with the same line $\mathbf Q_\ell v$. The universal
submodule $U_6$ is stable under $\operatorname{GSp}(V)$.
Its symplectic stability was proved in F.3.24; scalar matrices
act on degree-six derivations by their sixth power, so this
stability extends across the central isogeny
$\operatorname{Sp}(V)\times\mathbf G_m\to\operatorname{GSp}(V)$.

There is an explicit $\mathbf Q_\ell$-linear weight projector
\[
 e_{t^3}:U_6\cap T_{\mathrm{all}}
       \longrightarrow U_6\cap T_{\mathrm{all}}[t^3].          \tag{QCM.2}
\]
To construct it, choose $h\in T_{\mathrm{diag}}(\mathbf Q_\ell)$
at which $t^3$ differs from every other weight in the finite
dimensional derivation space. Such an $h$ exists because a
finite union of the proper zero sets of the Laurent
polynomials $\lambda-t^3$ does not cover the split torus
over the infinite field $\mathbf Q_\ell$. The Lagrange
interpolation polynomial in the operator $h$ gives (QCM.2).
Every factor preserves both displayed subspaces. The
equivariance of $q_6$ gives
\[
                              q_6 e_{t^3}=q_6.
\]
F.3.24 kills its target and proves (QCM.1).
In an arithmetic application the six tangency constraints
are first proved for an actual Galois element. The operator
(QCM.2) is then applied to that already established linear
constraint space inside the universal image. No lift of
this later algebraic projection to an arithmetic element
is asserted or needed. The same argument applies to the
four-line genus-two quotient calculation in F.3.22.



\paragraph{The local algebra and the nonsplit case.}
Put $B=J_C$ and $V=V_\ell B$. A degree-one arithmetic
normalizer line is stable on an open subgroup $G_L$.
Enlarge $L$ to define all geometric endomorphisms.
The semisimplicity and endomorphism statements of
Faltings, original Section~5, S\"atze~3--4, p.~360
\cite{GenFaltings1983}, give a rank-one idempotent
in $\operatorname{End}^0_{\Qb}(B)\otimes\Q_\ell$.
The exact rank formula and actual quotient construction
of F.3.18 apply. Since $B$ is absolutely simple, the
simple quotient there is isogenous to $B$ itself;
the formula forces the full geometric endomorphism
algebra to be a CM field $E$ of degree~6. The preceding
CM argument therefore makes $E/\Q$ Galois.

The comparison respecting actual endomorphisms makes
$V$ a free rank-one module over
\[
                E\otimes\Q_\ell=\prod_{\lambda\mid\ell}E_\lambda.
                                                               \tag{QCM.3}
\]
Every idempotent in this commutative product selects
some factors, and its $\Q_\ell$-rank is the sum of
the selected local degrees $[E_\lambda:\Q_\ell]$.
Galois transitivity on primes of $E$ makes all those
degrees equal, including at ramified primes.
Thus a rank-one idempotent exists only when every
local degree is one, precisely when $\ell$ splits
completely in $E$. At every other prime the
hypothesized degree-one normalizer line is excluded
by this actual endomorphism-algebra calculation.
This argument applies after every finite base
extension because the endomorphisms have already
all been defined.

\paragraph{The actual split-prime orbit of one normalizer.}
At a split prime the rank-one idempotent selects
exactly one embedding eigenline. Its actual
$G_{\Q}$ orbit is all six lines by (CMQ.7).
They are distinct as characters on every sufficiently
small open arithmetic subgroup: two equal characters
would allow a nonzero off-diagonal intertwiner,
contradicting Faltings' equality of the full commuting
algebra with the diagonal algebra (QCM.3).
The rational polarization supplies the three
nondegenerate conjugate pairs described above.

Choose a nonzero rational tangent vector at $y$.
It gives the actual tangential section of the
arithmetic group over $\Q$, as in F.3.24. For six
$\tau_i\in G_{\Q}$ carrying the initial eigenline
to the six specified lines, conjugate the original
procyclic subgroup and its normalizer by these
actual section lifts. Write the resulting subgroups
as $A_i$ and $N_i$. Their arithmetic images are
$\tau_i p(N)\tau_i^{-1}$ and are open. Choose a
finite Galois extension $L/\Q$ containing a field
defining all endomorphisms, with $G_L$ contained
in every one of these six images and in $G_K$.
Then each $N_i\cap p^{-1}(G_L)$ surjects onto
this same $G_L$. All groups act on the same
geometric punctured curve; no target curve or
homology space is replaced in this construction.

The arithmetic quotient and finite-base invariance
in Nakamura's original Lemma~4.5, Corollary~4.6,
and (4.7.1)--(4.7.2), pp.~115--117
\cite{AuditNakamura1996}, supply an actual
$\sigma\in G_L(6)$ whose degree-six class
$D\in U_6$ has $q_6(D)\ne0$. Here the quotient
is nonzero in degree~6 and has similitude character
$t^3$, with exactly the maps fixed in F.3.24.
For each $i$, lift this same $\sigma$ to $N_i$.
Because it acts trivially on $V$ and the linear
term of a generator of $A_i$ is nonzero, its
normalizer exponent is exactly one.

We spell out the rational correction so that no
primitivity of the original integral generator is
used. Scale that nonzero linear term to a basis
vector $v_i$, and write its full logarithm
$U_i=v_i+$ higher degrees. The canonical special
representative of $\sigma$ gives an equality
\[
             \widetilde\sigma(U_i)=\exp(\operatorname{ad}W_i)U_i
                                                               \tag{QCM.4}
\]
for a rational completed free-Lie element $W_i$.
Its degree-one term centralizes $v_i$, and hence
is a scalar multiple of $v_i$. Multiplying the
conjugating element by the appropriate rational
power of $\exp(U_i)$ removes that term while
leaving its action on $U_i$ unchanged. In each
successive degree below~6, the centralizer of
$v_i$ in the free Lie algebra forces the next
term to vanish. The degree-seven part of
(QCM.4) is therefore
\[
                         D(v_i)=[(W_i)_6,v_i].              \tag{QCM.5}
\]
This calculation is performed separately for each
$i$ with the same $D$. It proves that $D$ is
in $U_6\cap T_{\mathrm{all}}$. Equation (QCM.1)
then contradicts $q_6(D)\ne0$ and excludes the
split-prime degree-one case as well.

\paragraph{The exact passage from the first Magnus term to inertia.}
For completeness, the reduction used here is the
finite-cover argument of F.3.21 with the present
curve substituted through its actual maps. Let
$u$ generate the nontrivial procyclic subgroup.
A finite $\ell$-group quotient in which $u$ survives,
followed by the preimage of its cyclic image,
gives a finite etale cover $Y$ where $u$ has
nonzero rational abelian image. Its arithmetic
model exists over a finite extension. The
intersection of the normalizer with the covering
group still has open arithmetic image; its
geometric kernel acts inner on that covering
group and hence trivially on its homology.
Thus the detected homology line has exactly the
same character $\theta$ as the first nonzero
Magnus term $u_d\in\mathbb L_d(V)$.

The latter is pure of weight $d$ because the
original curve has exactly one puncture.
A nonzero image of the detected line in proper
homology of the compactification of $Y$ would,
at a common good Frobenius, force $d=1$.
That case was just excluded at every prime.
The line therefore lies in the actual kernel
of the homology map to the compactification,
namely the quotient of the boundary permutation
module tensored with $\Q_\ell(1)$ by its diagonal
copy. Over a finite extension defining all
boundary points, this entire kernel has character
$\chi_\ell$, so $\theta=\chi_\ell$ there.
The first Magnus term also shows that $\theta$
is trivial on the original geometric normalizer.
A lift from the intersected normalizer and any
lift from the full normalizer of the same
arithmetic element differ by that geometric
kernel. Hence the full normalizer over this
extension acts by $\chi_\ell$ and still has
open arithmetic image. Nakamura's original
Theorem~2.1.1, p.~94 \cite{AuditNakamura1994},
now places the original subgroup in actual
cusp inertia. Finite arithmetic base changes
leave the geometric group and this inertia
conjugacy class unchanged.
\end{proof}

\paragraph{Scope and source comparison.}
The general CM assertion has been derived from the exact
rational filtration and stabilizer maps, with no novelty
claim. For comparison, the opened Lemma~2.5 of
\cite{AuditGoodman2026} states the usual reflex-field
endomorphism formula under its explicit Galois-CM-field
hypothesis and points to Shimura's 1998 Proposition~30,
p.~65~\cite{AuditShimura1998Gate}; that original book passage has not been obtained
and is not an input to the proof above.
The curve theorem concerns a whole absolutely simple
Jacobian of dimension~3. A general detecting cover in
F.3.23 comes with its actual CM quotient, whose dimension
may be smaller than that of its Jacobian. The maps to
that quotient and its conjugates are retained; the
proof does not identify their orbit span with the whole
curve homology. In particular an original higher-genus
cover of $E_0\setminus O$ still has its norm quotient
to $E_0$. The full LM6 A.22 subgroup assertion remains
open in its stated scope.
